\section{Run the MMP on arithmetic threefolds}


\subsection{Pseudo-effective case}

    \begin{lemma}[Theorem 9.8]\label{lem:pe_terminal_modification_imply_special_termination}
        Let \((X,B)\) be a three-dimensional \(\bbQ\)-factorial dlt pair with \(\bbR\)-boundary which is projective and surjective over \(T\).
        Suppose that all three-dimensional \(\bbQ\)-factorial klt pairs \((X',B')\) with \(X'\) birational to \(X\) admit a terminal modification.
        Then for any sequence of steps in the \((K_X + B)\)-MMP, after finitely many steps, it will be isomorphism near \(\Supp B_i\).
    \end{lemma}

    \begin{lemma}\label{lem:terminal_modification}
        Let \((X,B)\) be a three-dimensional \(\bbQ\)-factorial klt pair with \(\bbR\)-boundary which is quasi projective and surjective over \(R\) and all residue fields has characteristic \(\neq 2,3,5\).
        Then there exists a projective birational morphism \(g:Y \to X\) and a terminal pair \((Y,B_Y)\) such that \(K_Y + B_Y = g^*(K_X + B)\).
    \end{lemma}
    \begin{proof}
        Take a resolution \(\pi:Y \to X\) with 
        \[ K_Y + B_Y + E = \pi^*(K_X + B) + F.\]
        We can assume that components of \(B_Y+E\) are disjoint.\Yang{Why?}
        Then \((Y,B_Y + E)\) is terminal.

        Run a \(K_Y + B_Y + E\)-MMP over \(X\), which is also a \(F\)-MMP.
        \Yang{By a difficulty argument, this MMP terminates.}
        Since \(F\) is exceptional over \(X\), by Negativity Lemma, \(F\) will be contracted in the MMP.
        Thus, the MMP terminates with a terminal model \((Y',B_{Y'}+E')\) such that \(K_{Y'} + B_{Y'} + E_{Y'} = \pi^*(K_X + B)\).
    \end{proof}

    \begin{theorem}\label{thm:run_mmp_pseudo_effective}
        Let \((X,B)\) be a \(\bbQ\)-factorial three-dimensional dlt pair with \(\bbR\)-boundary which is projective and surjective over \(T\) with \(\dim T >0\) and all residue fields has characteristic \(\neq 2,3,5\).
        Suppose that \(K_X + B\) is pseudo-effective.
        Then we can run a \((K_X + B)\)-MMP, and any sequence of steps in the MMP terminates.
        Consequently, we \((X,B)\) has a log minimal model.
    \end{theorem}
    \begin{proof}
        Let \(K_X + B \sim_T M \geq 0\).
        This can be checked by looking at the generic fiber and the non-vanishing theorem.

        We can run a \((K_X + B)\)-MMP over \(T\).

        After finitely many steps, it is isomorphic near \(\Supp \lfloor B \rfloor\) by Theorem 9.7.
        Hence it will be a \((K_X + \Delta)\)-MMP with \(\Delta = B - \lfloor B \rfloor + \varepsilon M\).
        By \cref{lem:pe_terminal_modification_imply_special_termination}, the MMP is isomorphic near \(\Supp B\) after finitely many steps.
        \Yang{Why isomorphic near \(M\)?}
        However, the extremal rays contracted are negativity on \(M\).
    \end{proof}


\subsection{Base point free theorem}

    \begin{proposition}\label{prop:abundance_slc_curve}
        Let \((X,B)\) be a slc curve with \(\bbQ\) boundary and \(K_X + B\) is nef.
        Then \(K_X + B\) is semi-ample.
    \end{proposition}
    \begin{proof}
        
    \end{proof}

    \begin{proposition}\label{prop:abundance_lc_surface}
        Let \((X,B)\) be a lc surface with \(\bbQ\)-boundary projective and surjective over \(T\) with \(\dim T >0\) and all residue fields has characteristic \(\neq 2,3,5\).
        Suppose that \(T = \Spec R\) with \((R,v,\kk)\) a local ring with residue field \(\kk\).
        Suppose that \(K_X + B\) is nef over \(T\).
        Then \(K_X + B\) is semi-ample over \(T\).
    \end{proposition}
    \begin{proof}
        It is enough to show \((K_X + B)|_{\Exc (K_X+B)}\) is semi-ample.

        By run MMP, we can assume that all components of \(\Exc (K_X+B)\) are either disjoint \(\lfloor B \rfloor\) or contained in \(\Supp B\).
        
        If disjoint, then linearly trivial.

        If contained in \(\Supp B\), then by adjunction and \cref{prop:abundance_slc_curve}.
    \end{proof}

    \begin{lemma}\label{lem:fiber_wise_base_point_free}
        Let \(X\) be a three-dimensional normal integral scheme which is projective and surjective over \(T\).
        Suppose that \(\calL\) is a nef but not big line bundle on \(X\) such that \(\calL|_{X_\eta}\) and \(\calL|_{X_{\bbQ}}\) are semi-ample, where \(\eta\) is the generic point of \(T\) and \(X_{\bbQ} = X \times_T \Spec \bbQ\).
        Then \(\calL\) is endowed with a map \(f: X \to V\) over \(T\) to an algebraic space \(V\) proper over \(T\).

        Moreover, if \(\calL|_{X_v}\) is semi-ample for every \(v \in V\), then \(\calL\) is semi-ample over \(T\).
    \end{lemma}
    \begin{proof}
        
    \end{proof}

    The idea is 
    \[ \text{abundance in slc curve } \implies \text{abundance in lc surface } \implies \text{bpf in klt threefold}. \]

    \begin{theorem}\label{thm:base_point_free}
        Let \((X,B)\) be a three-dimensional klt pair with \(\bbR\)-boundary which is projective and surjective over \(T\) with \(\dim T >0\) and all residue fields has characteristic \(\neq 2,3,5\).
        Suppose that \(\calL\) is a \(f\)-nef line bundle such that \(\calL - (K_X + B)\) is \(f\)-big and \(f\)-nef.
        Then \(\calL\) is \(f\)-semi-ample.
    \end{theorem}
    \begin{proof}
        By perturbing \(B\) slightly, we may assume that \(B\) is a \(\bbQ\)-divisor and \(L-(K_X+B)\) is \(f\)-ample.
        
        \begin{step}
            We can assume that we have a commutative diagram
            \[ \begin{tikzcd}
                X \arrow[rr, "g"] \arrow[dr, "f"] & & V \arrow[ld] \\
                & T & 
            \end{tikzcd} \]
            such that \(V = \Spec R\) with \((R,v,\kk)\) a local ring with residue field \(\kk\) of positive characteristic \(\neq 2,3,5\).
            And we just need to show that \(\calL|_{X_v}\) is \(f\)-semi-ample.
        \end{step}
        \begin{subproof}
            By previous results, we can assume that \(\calL\) is not big.
            Choose a general \(A \sim_{\bbQ} L - (K_X + B)\) such that \((X,B+A)\) is a klt pair.
            By bpf in lower dimension and \cref{lem:fiber_wise_base_point_free}, there is a map \(f: X \to V\) over \(T\) associated to \(L\) and we only need to check semi-ampleness of \(L\) fiber-wise over \(V\).
            Replacing \(V\) by its localization at \(v\) and \(X\) by the fiber product.
        \end{subproof}

        \begin{step}
            
        \end{step}


    \end{proof}

\subsection{Cone Theorem and Mori fiber space}

    \begin{lemma}[Corollary 9.18]\label{cor:contraction}
        Let \((X,B)\) be a \(\bbQ\)-factorial three-dimensional klt pair with \(\bbR\)-boundary which is projective and surjective over \(T\) with \(\dim T >0\) and all residue fields has characteristic \(\neq 2,3,5\).
        Suppose that \(\Sigma\) is a \((K_X+B)\)-negative extremal ray such that \(\Sigma = \overline{\symup{NE}}(X/T) \cap D^{\perp}\) with \(D\) nef and big.
        Then there exists a contraction \(f: X \to Z\) over \(T\) associated to \(\Sigma\).
    \end{lemma}

    \begin{theorem}\label{thm:cone_theorem}
        Let \((X,B)\) be a three-dimensional klt pair with \(\bbR\)-boundary which is projective and surjective over \(T\) with \(\dim T >0\) and all residue fields has characteristic \(\neq 2,3,5\).
        Then 
        \begin{enumerate}
            \item there are countably many rational curves \(C_i\) such that
                \[
                    \overline{NE}(X/T) = \overline{NE}(X/T)_{K_X + B \geq 0} + \sum_i \bbR_{\geq 0}[C_i];
                \]
            \item the rays \(\bbR_{\geq 0}[C_i]\) do not accumulate in the half-space \(N_1(X/T)_{K_X + B < 0}\);
            \item for each \(i\), we have 
                \[
                    -4d_{C_i} < -(K_X + B) \cdot C_i < 0.
                \]
        \end{enumerate}
    \end{theorem}

    \begin{lemma}\label{lem:bend_break_type_lemma}
        Let \((X,B)\) be a \(\bbQ\)-factorial klt threefold with \(\bbQ\)-boundary, projective and surjective over \(T\) with \(\dim T >0\) and all residue fields has characteristic \(\neq 2,3,5\).
        Suppose that \(K_X + B\) is not nef over \(T\).
        Then there exists \(n\) depending only on \((X,B)\) such that for any ample \(f\)-divisor \(H\), 
        \[ \lambda = \min \{t \geq 0 \mid K_X + B + tH \text{ is } f\text{ nef}\} \]
        is a rational number of form \(n/m\) for some \(m\).

        Furthermore, there exists a \(K_X+B+\lambda H\)-trivial curve \(\Gamma\) satisfying
        \[ 0 < -(K_X + B) \cdot \Gamma \leq 4d_\Gamma. \]
    \end{lemma}
    \begin{proof}
        Let \(L = K_X + B + \lambda H\).
        \begin{step}
            We show that \(\lambda \in \bbQ\) and there is a contraction \(f:X \to Z\) associated to \(L\) over \(T\).
        \end{step}
        \begin{subproof}
            If \(L\) is big, 
        \end{subproof}

        \begin{step}
            We can run a MMP.
        \end{step}

        \begin{step}
            Divisorial contraction case.
        \end{step}

        \begin{step}
            Flip case.
        \end{step}

        \begin{step}
            Mori fiber space case.
        \end{step}

        \begin{step}
            Finally, we show that \(n\) can be chosen uniformly.
        \end{step}
    \end{proof}

   